\chapter{Desarrollo}
\label{ch:desarrollo}

En este capítulo se detalla la resolución técnica del primer ejercicio práctico, correspondiente al diseño e
implementación del sistema de gestión de interrupciones externas utilizando la placa BluePill en un entorno Bare Metal.

\section{Ejercicio Práctico N1: interrupciones externas y el NVIC}
\label{sec.ej1:interrupciones}

En este ejercicio se ha configurado el microcontrolador STM32F103C8T6 para que el LED
interno (conectado al pin PC13) permanezca parpadeando periódicamente en el bucle principal, mientras que se
implementan dos pulsadores externos asociados a interrupciones externas con diferentes prioridades en el NVIC:
\begin{itemize}
    \item El Botón A (conectado en la GPIO C15) activa la interrupción de alta 
    prioridad (EXTI15\_10) que comanda el led externo 1 (conectado en la GPIO C14), haciéndolo parpadear 15 veces de forma rápida.
    \item El Botón B (conectado en la GPIO B0) activa la interrupción de baja prioridad (EXTI0)
    que comanda el led externo 2 (conectado en la GPIO B1), haciéndolo parpadear 15 veces de forma rápida.
\end{itemize}

El propósito de contar con este par de botones y leds es demostrar físicamente el mecanismo de la
anidación de interrupciones en el NVIC, observando cómo la interrupción de alta prioridad puede interrumpir la de
baja prioridad, deteniendo temporalmente su ciclo de parpadeo.

Para llevar a cabo la resolución de este ejercicio práctico, el proyecto se divide en tres archivos fundamentales: el
código de arranque en ensamblador (\texttt{crt.s}), el script de enlazado (\texttt{linker.ld}) y el código fuente
principal en C (\texttt{main.c}). A continuación, se detalla el análisis de cada uno de ellos y la relación funcional
que existe entre sí.

\subsection{Archivo de Inicialización y Tabla de Vectores}
\label{sec.ej1:inicializacion}

El archivo \texttt{crt.s}, cuyo código se expone en el listado~\ref{list.ej1:crt}, representa el código de inicio (startup)
del sistema. Al energizar el microcontrolador, la CPU del Cortex-M3 inicializa el puntero de pila (Stack Pointer) con el
valor ubicado en la dirección de memoria \texttt{0x08000000} (inicio de la memoria FLASH) y salta a la dirección del vector
de reinicio (Reset Vector) ubicada en \texttt{0x08000004}.

En este archivo se define la sección \texttt{.isr\_vector}, la cual contiene la tabla de vectores de interrupción requerida por
la arquitectura ARM Cortex-M3~\cite{cortexm3trm}. Las primeras posiciones corresponden a excepciones internas del núcleo (Reset, NMI y Hard Fault).
Las posiciones del núcleo restantes (de la 4 a la 15) se rellenan con ceros mediante la directiva \texttt{.rept 12}, ya que no
son utilizadas en esta aplicación.

A continuación de las excepciones del núcleo, se posicionan los vectores de interrupción específicos de los
periféricos de la placa STM32. Para demostrar el manejo de prioridades e interrupción anidada, se implementan
dos controladores de interrupción diferentes: \texttt{EXTI0\_IRQHandler} (asociado a la línea EXTI0, posición
física 6 de los periféricos y 22 global) y \texttt{EXTI15\_10\_IRQHandler} (asociado a las líneas EXTI15 a EXTI10,
posición física 40 de los periféricos y 56 global). Para ubicarlos de forma exacta en la FLASH, se rellenan
las posiciones intermedias con ceros usando las directivas \texttt{.rept 6} y \texttt{.rept 33}.

El punto de entrada físico del firmware es el manejador de reset \texttt{\_reset}, el cual llama a la función principal en C
mediante la instrucción de ramificación con enlace \texttt{bl main}. Si la función principal retornara, el microcontrolador
quedaría en un bucle infinito para evitar comportamientos erráticos.

\lstinputlisting[language={ARM-Assembly}, caption={Archivo de inicialización y tabla de vectores (\texttt{crt.s}).}, label={list.ej1:crt}]{sources/src/crt.s}

\subsection{Script de Enlazado}
\label{sec.ej1:enlazador}

El archivo \texttt{linker.ld}, presentado en el listado~\ref{list.ej1:linker}, es el script de configuración del enlazador
(GNU Linker). Su función principal es indicar cómo deben distribuirse las secciones de código y datos dentro del espacio de
direcciones de memoria del microcontrolador.

En primer lugar, define el bloque \texttt{MEMORY}, especificando que el chip STM32F103C8T6 posee una memoria FLASH de lectura
y ejecución (\texttt{rx}) de $64\kB$ con origen en \texttt{0x08000000}, y una memoria RAM de lectura, escritura y ejecución
(\texttt{xrw}) de $20\kB$ con origen en \texttt{0x20000000}.

El símbolo \texttt{\_esstack} se calcula sumando el origen de la RAM con su longitud, apuntando así a la última posición de la
memoria RAM (pila invertida o descendente). Este símbolo es exportado y consumido directamente en la primera posición de la
tabla de vectores en \texttt{crt.s}.

Dentro del bloque \texttt{SECTIONS}, la directiva \texttt{KEEP(*(.isr\_vector))} asegura que la sección de la tabla de vectores
de interrupción no sea eliminada por optimizaciones del compilador y que se ubique de manera obligatoria al principio de la
memoria FLASH. Posteriormente, se disponen las secciones de código y constantes (\texttt{.text} y \texttt{.rodata}) en la FLASH,
mientras que las secciones de variables globales (\texttt{.data} y \texttt{.bss}) se mapean para su uso en la RAM.

\lstinputlisting[language={LinkerScript}, caption={Script de enlazado de memoria (\texttt{linker.ld}).}, label={list.ej1:linker}]{sources/src/linker.ld}

\subsection{Código Fuente Principal en C}
\label{sec.ej1:codigo_c}

El archivo \texttt{main.c}, expuesto en el listado~\ref{list.ej1:main}, implementa la inicialización de los registros
de control y la rutina de servicio para procesar las interrupciones externas. Al no contar con los archivos de cabecera estándar del fabricante,
se definen de manera explícita las direcciones físicas de los periféricos basándose en el manual de referencia~\cite{stm32f103rm}.

\subsubsection{Mapeo de Registros}

\begin{itemize}
    \item \texttt{GPIOB}: Con dirección base en \texttt{0x40010C00}, define los registros \texttt{GPIOB\_CRL}
    y \texttt{GPIOB\_ODR} para configurar el segundo botón (PB0) como entrada con pull-up y el segundo led
    externo (PB1) como salida digital push-pull.
    \item \texttt{GPIOC}: Dirección base en \texttt{0x40011000}, define el registro de configuración alto
    \texttt{GPIOC\_CRH} para los pines 13, 14 y 15, \texttt{GPIOC\_ODR} para el manejo de los estados de
    los pines de salida, y \texttt{GPIOC\_BSRR} para la alternancia atómica del led indicador interno.
    \item \texttt{RCC}: Dirección base \texttt{0x40021000}. Se utiliza \texttt{RCC\_APB2ENR} para habilitar
    los relojes de los periféricos GPIOB, GPIOC y AFIO.
    \item \texttt{AFIO}: Dirección base \texttt{0x40010000}. Se incorporan los registros \texttt{AFIO\_EXTICR1}
    y \texttt{AFIO\_EXTICR4} para asociar los pines PB0 y PC15 a los canales EXTI0 y EXTI15 respectivamente.
    \item \texttt{EXTI}: Dirección base \texttt{0x40010400}. Se emplean los registros de máscara de interrupción
    (\texttt{EXTI\_IMR}), detección por flanco de bajada (\texttt{EXTI\_FTSR}) y el registro de petición pendiente
    (\texttt{EXTI\_PR}) para limpiar el flag de interrupción.
    \item \texttt{NVIC}: Se emplean los registros \texttt{NVIC\_ISER0} y \texttt{NVIC\_ISER1} para habilitar los
    canales EXTI0 (IRQ 6) y EXTI15\_10 (IRQ 40), y los registros de prioridad \texttt{NVIC\_IPR1} y
    \texttt{NVIC\_IPR10} para configurar sus prioridades relativas.
\end{itemize}

\subsubsection{Configuración Paso a Paso en la Función Principal}

Dentro del ciclo de inicialización en la función \texttt{main()}, se ejecutan las siguientes acciones en secuencia:
\begin{enumerate}
    \item Habilitación del reloj de periféricos: Se configuran los bits 4, 3 y 0 de \texttt{RCC\_APB2ENR} para
    suministrar señal de reloj a los puertos GPIOC, GPIOB y al módulo AFIO.
    \item Configuración de PC13, PC14 y PB1 como salidas: Se establecen en modo push-pull de $50\MHz$.
    PC13 y PC14 se configuran en el puerto C a través de \texttt{GPIOC\_CRH}, mientras que PB1 se
    configura en el puerto B usando \texttt{GPIOB\_CRL}. PC13 controla el led interno, PC14 el led externo 1,
    y PB1 el led externo 2.
    \item Configuración de entradas con Pull-up: Se configura PC15 (Botón A) y PB0 (Botón B) en modo entrada con
    resistencia interna escribiendo \texttt{0x8} en sus respectivos registros de configuración y estableciendo
    en '1' los bits correspondientes en los registros de salida de datos (\texttt{GPIOx\_ODR}) para definir
    las pull-up.
    \item Asociación de interrupciones externas: Se escribe \texttt{0x2} en el campo EXTI15 de \texttt{AFIO\_EXTICR4}
    y \texttt{0x1} en el campo EXTI0 de \texttt{AFIO\_EXTICR1} para direccionar los pines PC15 y PB0 a las líneas
    de interrupción externa correspondientes.
    \item Configuración del disparo y máscara: Se configuran los flancos descendentes escribiendo un '1' en los
    bits 15 y 0 de \texttt{EXTI\_FTSR}, y se habilitan las líneas escribiendo en \texttt{EXTI\_IMR}.
    \item Configuración de prioridades de interrupción: Para demostrar la anidación, se le asigna prioridad alta
    ($0x40$) a EXTI15\_10 mediante el registro \texttt{NVIC\_IPR10}, y una prioridad menor ($0x80$) a EXTI0 a través
    del registro \texttt{NVIC\_IPR1}.
    \item Habilitación en el NVIC: Se escribe un '1' en el bit 6 de \texttt{NVIC\_ISER0} (para activar EXTI0) y en el
    bit 8 de \texttt{NVIC\_ISER1} (para activar EXTI15\_10).
\end{enumerate}

Una vez completada la inicialización, la CPU ejecuta de forma ininterrumpida el bucle principal \texttt{while(1)},
que alterna el estado del LED interno en PC13 mediante retardos por software (bucles vacíos de larga duración).

\subsubsection{Rutina de Servicio de Interrupción (ISR) y Demostración de Prioridades}

La demostración práctica del manejo de prioridades e interrupción anidada (preemption) se ejecuta al interactuar
con las rutinas de servicio definidas en \texttt{main.c}:

\begin{itemize}
    \item \texttt{EXTI15\_10\_IRQHandler} (Prioridad Alta - 0x40): Al presionar el Botón A en PC15, se inicia un bucle
    largo observable por hardware que hace parpadear el led externo 1 (PC14) rápidamente 15 veces seguidas mediante
    un retraso por software.
    \item \texttt{EXTI0\_IRQHandler} (Prioridad Baja - 0x80): Al presionar el Botón B en PB0, se inicia un bucle largo
    similar que hace parpadear el led externo 2 (PB1) de manera rápida 15 veces seguidas.
\end{itemize}

El comportamiento observable del sistema valida de forma empírica la jerarquía de prioridades del NVIC:
\begin{enumerate}
    \item Si se presiona el Botón B (baja prioridad), el led 2 en PB1 comienza a parpadear rápidamente.
    \item Si durante este parpadeo se presiona el Botón A (alta prioridad), el NVIC suspende de inmediato la ejecución
    de la ISR de menor prioridad y salta a \texttt{EXTI15\_10\_IRQHandler}. Esto se evidencia físicamente porque el
    led 1 en PC14 comienza a parpadear de inmediato, pausando el parpadeo del led 2.
    \item Al finalizar el parpadeo del led 1, el procesador retorna a la rutina de menor prioridad, reanudando
    el parpadeo del led 2 en el punto exacto donde fue pausado.
    \item Por el contrario, si se presiona primero el Botón A, el led 1 parpadea y cualquier pulsación del Botón B
    no tiene efecto inmediato; la ISR de menor prioridad queda retenida en estado pendiente y solo comienza a
    parpadear el led 2 una vez que el led 1 ha completado sus 15 destellos.
\end{enumerate}

\lstinputlisting[language={C}, caption={Código fuente principal e implementación de la ISR (\texttt{main.c}).}, label={list.ej1:main}]{sources/src/main.c}

\subsection{Relación y Flujo de Ejecución entre los Componentes}
\label{sec.ej1:relacion}

La interacción de los tres archivos de desarrollo analizados conforma la base necesaria para la correcta ejecución
en Bare Metal:
\begin{enumerate}
    \item El archivo de enlazado (\texttt{linker.ld}) asigna de manera rígida los límites de la memoria física y
    establece que la tabla de vectores definida en \texttt{crt.s} comience exactamente en la posición \texttt{0x08000000}.
    \item El código en ensamblador (\texttt{crt.s}) posiciona la dirección física de la rutina
    \texttt{EXTI15\_10\_}\allowbreak\texttt{IRQHandler} en el vector número 40 del NVIC (posición 56 de la tabla).
    \item Durante la compilación y el enlace, el compilador traduce la función definida en \texttt{main.c} en una
    dirección física de memoria, la cual el enlazador copia en el vector asignado dentro de la tabla de interrupciones.
    \item Cuando la interrupción física se produce en el microcontrolador, el NVIC consulta la dirección de memoria
    asociada al vector 40 en la FLASH y salta hacia la función implementada en C, logrando la interacción asíncrona.
\end{enumerate}
